\section*{Capítulo \ref{ch:b2:acyl-substitution} --- Derivados de los ácidos carboxílicos}

\begin{solution}{exo:b2:acyl-substitution:1}
Etanamida $<$ etanoato de etilo $<$ anhídrido etanoico $<$ cloruro de etanoílo.
\end{solution}

\begin{solution}{exo:b2:acyl-substitution:2}
Ácido etanoico (+ \ce{HCl}); etanoato de etilo; etanamida (+ cloruro de amonio: el segundo equivalente
de amoníaco atrapa el \ce{HCl}); anhídrido etanoico (+ \ce{NaCl});
\textit{N},\textit{N}-dimetiletanamida (+ cloruro de trietilamonio).
\end{solution}

\begin{solution}{exo:b2:acyl-substitution:3}
\ce{C6H5COOCH3 + NaOH -> C6H5COONa + CH3OH}. $13.6/136.15 = \qty{0.0999}{mol}$, y otro tanto de
\ce{NaOH}: \qty{4.00}{g}.
\end{solution}

\begin{solution}{exo:b2:acyl-substitution:4}
$\gamma$-Butirolactona (oxolan-2-ona): un anillo de cinco eslabones, cuatro carbonos y un oxígeno.
\end{solution}

\begin{solution}{exo:b2:acyl-substitution:5}
El amoníaco se adiciona al carbono carbonílico; un protón pasa del nitrógeno al oxígeno (o al
grupo saliente); el etóxido (como etanol, tras una transferencia de protón) sale y el \ce{C=O} se forma de nuevo:
etanamida y etanol.
\end{solution}

\begin{solution}{exo:b2:acyl-substitution:6}
Se acila el \ce{OH} fenólico (catálisis ácida): ácido 2-acetoxibenzoico; subproducto, el ácido etanoico.
\end{solution}

\begin{solution}{exo:b2:acyl-substitution:7}
2-Metilpropan-2-ol. La primera adición da propanona, más electrófila que el éster, que
toma de inmediato el segundo metilo.
\end{solution}

\begin{solution}{exo:b2:acyl-substitution:8}
Butano-1,4-diol: el éster se reduce a dos alcoholes primarios, y el anillo se abre.
\end{solution}

\begin{solution}{exo:b2:acyl-substitution:9}
Un equivalente: $x^2/(1 - x)^2 = 4$, $x = 2/3$. Tres equivalentes: $x^2/[(1 - x)(3 - x)] = 4$, $3x^2 - 16x
+ 12 = 0$, $x = 0.90$. Eliminar el agua a medida que se forma (Dean--Stark) la desplaza aún más.
\end{solution}

\begin{solution}{exo:b2:acyl-substitution:10}
\ce{SOCl2} da cloruro de benzoílo; con dos equivalentes de dietilamina (o uno más trietilamina)
da la amida. El ácido y la amina mezclados directamente sufren una reacción ácido--base: benzoato de
dietilamonio, que solo da la amida calentando fuertemente.
\end{solution}

\begin{solution}{exo:b2:acyl-substitution:11}
Tres KOH por trioleína: $3 \times 56.11/885.4 = \qty{0.190}{g}$ por gramo, un índice de \omterm{def:b2:acyl-substitution:saponification}{saponificación} de
190.
\end{solution}

\begin{solution}{exo:b2:acyl-substitution:12}
\ce{NaBH4}: solo la cetona, a alcohol secundario. \ce{LiAlH4}: la cetona al alcohol
secundario, el éster a un alcohol primario (liberando la parte alcohólica), la amida a una amina.
\end{solution}

\begin{solution}{pb:b2:acyl-substitution:1}
\textbf{1.} Ácido oleico \ce{C18H34O2}; glicerol \ce{C3H8O3}.
\textbf{2.} \ce{C57H104O6}: $57 \times 12.011 + 104 \times 1.008 + 6 \times 15.999 = \qty{885.45}{g/mol}$.
\textbf{3.} Tres.
\textbf{4.} El doble enlace cis acoda las cadenas, que se empaquetan mal: fuerzas intermoleculares débiles, un
punto de fusión bajo.
\textbf{5.} Glicerol y tres oleatos de sodio, un jabón.
\textbf{6.} Una mezcla de ésteres metílicos de ácidos grasos.
\textbf{7.} \ce{C57H104O6 + 3CH3OH -> C3H8O3 + 3C19H36O2}.
\textbf{8.} El ion metóxido, formado a partir del metanol y la base, es un nucleófilo mucho más fuerte que el
metanol; ataca los carbonilos de los ésteres.
\textbf{9.} El carbono carbonílico de un grupo éster, que lleva \ce{O-}, el \ce{OCH3} del metóxido y el
oxígeno del glicerilo: tetraédrico.
\textbf{10.} La \omterm{def:b2:acyl-substitution:transesterification}{transesterificación} es un equilibrio con una constante cercana a 1; el exceso de metanol la desplaza
hacia los ésteres metílicos.
\textbf{11.} El glicerol se separa como capa inferior: retirar un producto desplaza el equilibrio.
\textbf{12.} No: el metóxido se regenera en cada intercambio.
\textbf{13.} Neutraliza la base, dando oleato de sodio (jabón): el catalizador se consume.
\textbf{14.} Los saponifica: el hidróxido (procedente del agua y del metóxido) corta los ésteres en jabón y
metanol, de forma irreversible.
\textbf{15.} Consume base, emulsiona la mezcla e impide que la capa de glicerol se separe.
\textbf{16.} Mediante una \omterm{def:b2:acyl-substitution:esterification}{esterificación} previa con metanol catalizada por ácido, que convierte los ácidos libres en
ésteres metílicos.
\textbf{17.} El agua hidroliza los ésteres y convierte el catalizador en hidróxido, que saponifica.
\textbf{18.} $10^6/885.45 = \qty{1129}{mol}$.
\textbf{19.} $3 \times 1129 \times 32.04 = \qty{108.6}{kg}$.
\textbf{20.} $3 \times 1129 \times 296.50 = \qty{1004.6}{kg}$.
\textbf{21.} Entrada: $1000 + 108.6 = \qty{1108.6}{kg}$; salida: $104.0 + 1004.6 = \qty{1108.6}{kg}$.
\textbf{22.} \qty{20}{kg} de ácido oleico, $20\,000/282.47 = \qty{70.8}{mol}$, y otro tanto de \ce{NaOH}:
\qty{2.83}{kg}.
\textbf{23.} $1129 \times 92.09 = \boldsymbol{\approx \qty{104}{kg}}$ de glicerol por tonelada.
\end{solution}
